% Source-derived chapter 03: Analytic side
% Original source: anticyclotomic-iwasawa-rational-elliptic-eisenstein
\chapter{Analytic side}
\label{anticyclotomic-iwasawa-rational-elliptic-eisenstein-chapter-03}
\source{anticyclotomic-iwasawa-rational-elliptic-eisenstein}

\begin{remark}[Source section]
\label{anticyclotomic-iwasawa-rational-elliptic-eisenstein-chapter-03-source}
\source{anticyclotomic-iwasawa-rational-elliptic-eisenstein}
\source{anticyclotomic-iwasawa-rational-elliptic-eisenstein}
This chapter reproduces the corresponding section of the retrieved
LaTeX source, including its definitions, statements, examples, and
proofs. It has not yet been translated into Lean.
\notready
\end{remark}

% The source section title was: Analytic side
\label{IMCII}
\source{anticyclotomic-iwasawa-rational-elliptic-eisenstein}

Let $E/\Q$ be an elliptic curve of conductor $N$, $p\nmid 2N$ a prime of good reduction for $E$, and $K$ an imaginary quadratic field satisfying hypotheses (\ref{eq:intro-Heeg}), (\ref{eq:intro-spl}), and (\ref{eq:intro-disc}) from the introduction; in particular, $p=v\bar{v}$ splits in $K$.

In this section, assuming $E[p]^{ss}=\mathbb{F}_p(\phi)\oplus\mathbb{F}_p(\psi)$ as $G_\Q$-modules, we prove an analogue of  Theorem~\ref{MAINalgside} on the analytic side, relating the Iwasawa invariants of an anticyclotomic $p$-adic $L$-function of $E$ to the Iwasawa invariants of anticyclotomic Katz $p$-adic $L$-functions attached to $\phi$ and $\psi$. 


\subsection{\texorpdfstring{$p$}{p}-adic \texorpdfstring{$L$}{L}-functions}\label{subsec:Lp}
\source{anticyclotomic-iwasawa-rational-elliptic-eisenstein}

Recall that $\Lambda=\Z_p\llbracket\Gamma\rrbracket$ denotes the anticyclotomic Iwasawa algebra, and set $\Lambda^{\rm ur}=\Lambda\hat{\otimes}_{\Z_p}\Z_p^{\rm ur}$,  
for $\Z_p^{\rm ur}$ the completion of the ring of integers of the maximal unramified extension of $\Q_p$. 

We shall say that an algebraic Hecke character $\psi:K^\times\backslash\mathbb{A}_K^\times\rightarrow\mathbb{C}^\times$ has infinity type $(m,n)$ if 
the component $\psi_\infty$ of $\psi$ at $\infty$ satisfies
$\psi_\infty(z)=z^{m}\bar{z}^{n}$ for all $z\in (K\otimes\mathbb{R})^\times\simeq\mathbb{C}^\times$, 
where the last identification is made via $\iota_\infty$.


\subsubsection{The Bertolini--Darmon--Prasanna \texorpdfstring{$p$}{p}-adic \texorpdfstring{$L$}{l}-functions} 

Fix an integral ideal $\mathfrak{N}\subset\mathcal{O}_K$ with 
\begin{equation}\label{eq:N}
\source{anticyclotomic-iwasawa-rational-elliptic-eisenstein}
\mathcal{O}_K/\mathfrak{N}\simeq\Z/N\Z.
\end{equation}
Let $f\in S_2(\Gamma_0(N))$ be the newform associated with $E$. Following \cite{BDP}, one has the following result. 

\begin{theorem}
\source{anticyclotomic-iwasawa-rational-elliptic-eisenstein}
\notready\label{thm:BDP}
\source{anticyclotomic-iwasawa-rational-elliptic-eisenstein}
There exists an element $\Lcal_E\in\Lambda^{\rm ur}$ characterized by the following interpolation property: For every character $\xi$ of $\Gamma$ crystalline at both $v$ and $\bar{v}$ and corresponding to a Hecke character of $K$ of infinity type $(n,-n)$ with $n\in\Z_{>0}$ and $n\equiv 0\pmod{p-1}$, we have
\[
\mathcal{L}_E(\xi)=\frac{\Omega_p^{4n}}{\Omega_\infty^{4n}}\cdot\frac{\Gamma(n)\Gamma(n+1)\xi(\mathfrak{N}^{-1})}{4(2\pi)^{2n+1}\sqrt{D_K}^{2n-1}}\cdot\bigl(1-a_p\xi(\bar{v})p^{-1}+\xi(\bar{v})^2p^{-1}\bigr)^2\cdot L(f/K,\xi,1),
\]  
where $\Omega_p$ and $\Omega_\infty$ are CM periods attached to $K$ as in \cite[\S{2.5}]{cas-hsieh1}.
\end{theorem}

\begin{proof}
This was originally constructed in \cite{BDP} as a continuous function of $\xi$, and later explicitly constructed as a measure in \cite{cas-hsieh1} (following the approach in \cite{brakocevic}). Since this refined construction will be important for our purposes in this section, we recall some of the details. 

Let ${\rm Ig}(N)$ be the Igusa scheme over $\Z_{(p)}$ parametrizing elliptic curves with $\Gamma_1(Np^\infty)$-level structure as in \cite[\S{2.1}]{cas-hsieh1};  its complex points admit a uniformization 
\begin{equation}\label{eq:unif}
\source{anticyclotomic-iwasawa-rational-elliptic-eisenstein}
[\;,\;]:\mathfrak{H}\times{\rm GL}_2(\hat{\Q})\rightarrow{\rm Ig}(N)(\mathbb{C}).
\end{equation} 
Let $c$ be a positive integer prime to $Np$. Then $\vartheta:=(D_K+\sqrt{-D_K})/2$ and the element $\xi_c:=\varsigma^{(\infty)}\gamma_c\in\mathrm{GL}_2(\hat{\mathbb{Q}}^{})$ constructed in \cite[p.~577]{cas-hsieh1} define a point 
\[
x_c:=[(\vartheta,\xi_c)]\in\mathrm{Ig}(N)(\mathbb{C})
\]
rational over $K[c](v^\infty)$, the compositum of the ring class field $K$ of conductor $c$ and the ray class field of $K$ of conductor $v^\infty$. For every $\mathcal{O}_c$-ideal $\mathfrak{a}$ prime to $\mathfrak{N}v$, let $a\in\hat{K}^{(cp),\times}$ be such that $\mathfrak{a}=a\hat{\mathcal{O}}_c\cap K$ and set
\[
\sigma_\mathfrak{a}:=\mathrm{rec}_K(a^{-1})\vert_{K[c](v^\infty)}\in\mathrm{Gal}(K[c](v^\infty)/K),
\]
where ${\rm rec}_K:K^\times\backslash\hat{K}^\times\rightarrow G_K^{\rm ab}$ is the reciprocity map (geometrically normalized). Then by Shimura's reciprocity law, the point $x_\mathfrak{a}:=x_c^{\sigma_\mathfrak{a}}$ is defined by the pair $(\vartheta,\overline{a}^{-1}\xi_c)$ under $(\ref{eq:unif})$.  

Let $V_p(N;R)$ be the space of $p$-adic modular forms of tame level $N$ defined over a $p$-adic ring $R$ (as recalled in \cite[\S{2.2}]{cas-hsieh1}), and let $S_{\mathfrak{a}}\hookrightarrow{\rm Ig}(N)_{/\Z_p^{\rm ur}}$ be the local deformation space of $x_\mathfrak{a}\otimes\overline{\mathbb{F}}_p\in{\rm Ig}(N)(\overline{\mathbb{F}}_p)$, so we have $\mathcal{O}_{S_{\mathfrak{a}}}\simeq\Z_p^{\rm ur}\llbracket t_\mathfrak{a}-1\rrbracket$ by Serre--Tate theory. Viewing $f$ as a $p$-adic modular form, the Serre--Tate expansion 
\[
f(t_{\mathfrak{a}}):=f\vert_{S_\mathfrak{a}}\in\bZ_p^{\rm ur}\llbracket t_{\mathfrak{a}}-1\rrbracket 
\]
defines a $\Z_p^{\rm ur}$-valued measure ${\rm d}\mu_{f,\mathfrak{a}}$ on $\Z_p$ characterized (by Mahler's theorem, see e.g. \cite[\S{3.3}, Thm.~1]{hida-blue}) by
\begin{equation}\label{eq:def-meas}
\source{anticyclotomic-iwasawa-rational-elliptic-eisenstein}
\int_{\Z_p}\binom{x}{m}{\rm d}\mu_{f,{\mathfrak{a}}}=\binom{\theta}{m}f(x_{\mathfrak{a}})
\end{equation}
for all $m\geqslant 0$, %(see Lemma~3.1), 
where $\theta:V_p(N;\Z_p)\rightarrow V_p(N;\Z_p)$ is the Atkin--Serre operator, acting as $qd/dq$ on $q$-expansions. Similarly, the $p$-depletion 
\[
f^\flat=\sum_{p\nmid n}a_nq^n
\]
defines a $\Z_p^{\rm ur}$-valued measure ${\rm d}\mu_{f^\flat,\mathfrak{a}}$  on $\Z_p$ (supported on $\Z_p^\times$) with $p$-adic Mellin transform $f^\flat(t_\mathfrak{a})$, and we let ${\rm d}\mu_{f^\flat_{\mathfrak{a}}}$ be the measure on $\Z_p^\times$ corresponding to $f(t_{\mathfrak{a}}^{{\rm N}(\mathfrak{a})^{-1}\sqrt{-D_K}^{-1}})$ (see \cite[Prop.~3.3]{cas-hsieh1}).

Letting $\eta$ be an auxiliary anticyclotomic Hecke character of $K$ of infinity type $(1,-1)$ and conductor $c$, define $\mathscr{L}_{v,\eta}\in\Lambda^{\mathrm{ur}}$ by
\begin{equation}\label{eq:cusp-measure}
\source{anticyclotomic-iwasawa-rational-elliptic-eisenstein}
\mathscr{L}_{v,\eta}(\phi)=
\sum_{[\mathfrak{a}]\in\mathrm{Pic}(\mathcal{O}_c)}
\eta(\mathfrak{a})\mathbf{N}(\mathfrak{a})^{-1}
\int_{\Z_p^\times}\eta_v(\phi\vert[\mathfrak{a}])\;\mathrm{d}\mu_{f^\flat_\mathfrak{a}}
\end{equation}
where $\eta_v$ denotes the $v$-component of $\eta$, and $\phi\vert[\mathfrak{a}]:\Z_p^\times\rightarrow\mathcal{O}_{\mathbb{C}_p}^\times$ is defined by $(\phi\vert[\mathfrak{a}])(x)=\phi({\rm rec}_v(x)\sigma_{\mathfrak{a}}^{-1})$ for the local reciprocity map ${\rm rec}_v:K_v^\times\rightarrow G_K^{\rm ab}\twoheadrightarrow\Gamma$. Then by \cite[Prop.~3.8]{cas-hsieh1} the element $\mathcal{L}_E\in\Lambda^{\rm ur}$
defined by
\[
\mathcal{L}_E(\xi):=\mathscr{L}_{v,\eta}(\eta^{-1}\xi)^2
\]
has the stated interpolation property. 
\end{proof}


\subsubsection{Katz $p$-adic L-functions} 

Let $\theta:G_\Q\rightarrow\mathbb{Z}_p^\times$ be a Dirichlet character of conductor $C$. As it will suffice for our purposes, we assume that $C\vert N$ (so $p\nmid C$), and let $\mathfrak{C}\vert\mathfrak{N}$ be such that $\mathcal{O}_K/\mathfrak{C}=\Z/C\Z$. 

The next result follows from the work of Katz \cite{katz}, as extended by Hida--Tilouine \cite{hidatilouineI}.

\begin{theorem}
\source{anticyclotomic-iwasawa-rational-elliptic-eisenstein}
\notready\label{thm:Katz}
\source{anticyclotomic-iwasawa-rational-elliptic-eisenstein}
There exists an element $\mathcal{L}_{\theta}\in\Lambda^{\rm ur}$ characterized by the following interpolation property: For every character $\xi$ of $\Gamma$ crystalline at both $v$ and $\bar{v}$ and corresponding to a Hecke character of $K$ of infinity type $(n,-n)$ with $n\in\Z_{>0}$ and $n\equiv 0\pmod{p-1}$, we have
\begin{align*}
\Lcal_{\theta}(\xi)=\frac{\Omega_p^{2n}}{\Omega_\infty^{2n}}\cdot 4\Gamma(n+1)\cdot\frac{(2\pi i)^{n-1}}{\sqrt{D_K}^{n-1}}&\cdot\bigl(1-\theta^{-1}(p)\xi^{-1}(v)\bigr)\cdot\bigl(1-\theta(p)\xi(\bar{v})p^{-1})\bigr)\\
&\times\prod_{\ell\vert C}(1-\theta(\ell)\xi(w)\ell^{-1})\cdot L(\theta_{K}\xi\mathbf{N}_K,0),
\end{align*}
where $\Omega_p$ and $\Omega_\infty$ are as in Theorem~\ref{thm:BDP}, and for each $\ell\vert C$ we take the prime $w\vert\ell$ with $w\vert\mathfrak{C}$.
\end{theorem}

\begin{proof}
The character $\theta$  (viewed as a character of $K$) defines a projection 
\[
\pi_\theta:\bZ_p^{\rm ur}\llbracket{\rm Gal}(K(\mathfrak{C}p^\infty)/K)\rrbracket\rightarrow\Lambda^{\rm ur},
\]
where $K(\mathfrak{C}p^\infty)$ is the ray class field of $K$ of conductor $\mathfrak{C}p^\infty$
(this projection is just $g\mapsto \theta(g)[g]$ for $g\in {\rm Gal}(K(\mathfrak{C}p^\infty)/K$ 
and $[g]$ the image of $g$ in $\Gamma$).
The element $\Lcal_{\theta}$ is then obtained by applying $\pi_\theta$ to the Katz $p$-adic $L$-function described in \cite[Thm.~27]{kriz}, setting $\chi^{-1}=\theta_K\xi\mathbf{N}_K$. 
\end{proof}


\subsection{Comparison II: Analytic Iwasawa invariants}


The following theorem follows from the main result of \cite{kriz}. Following the notations in \emph{op.\,cit},
we let $N_0$ be the square-full part of $N$ (so the quotient $N/N_0$ is square-free), and fix an integral ideal $\mathfrak{N}\subset\mathcal{O}_K$ as in (\ref{eq:N}).

Let also $f=\sum_{n=1}^\infty a_nq^n\in S_2(\Gamma_0(N))$ be the newform associated with $E$, and denote by $\lambda^\iota$ the image of $\lambda\in\Lambda$ under the involution of $\Lambda$ given by $\gamma\mapsto\gamma^{-1}$ for $\gamma\in\Gamma$. 


\begin{theorem}
\source{anticyclotomic-iwasawa-rational-elliptic-eisenstein}
\notready\label{thm:kriz}
\source{anticyclotomic-iwasawa-rational-elliptic-eisenstein}
Assume that $E[p]^{ss}\simeq\mathbb{F}_p(\phi)\oplus\mathbb{F}_p(\psi)$ as $G_\bQ$-modules,
with the characters $\phi, \psi$ labeled so that $p\nmid\mathrm{cond}(\phi)$, and suppose $\phi\neq\mathds{1}$.
Then there is a factorization $N/N_0=N_+N_-$ with
\[
%a_\ell\equiv
\begin{cases}
a_\ell\equiv\phi(\ell)\pmod{p}&\textrm{if $\ell\vert N_+$,}\\
a_\ell\equiv\psi(\ell)\pmod{p}&\textrm{if $\ell\vert N_-$,}\\
a_\ell\equiv 0\pmod{p}&\textrm{if $\ell\vert N_0,$}
\end{cases}
\]
such that the following congruence holds
\[
\mathcal{L}_E\equiv(\mathcal{E}_{\phi,\psi}^\iota)^2\cdot(\mathcal{L}_\phi)^2\pmod{p\Lambda^{\rm ur}},
\]
where 
\[
\mathcal{E}_{\phi,\psi}=\prod_{\ell\vert N_0N_-}\mathcal{P}_w(\phi)\cdot\prod_{\ell\vert N_0N_+}\mathcal{P}_w(\psi),
\] 
and for each $\ell\vert N$ we take the prime $w\vert\ell$ with $w\vert\mathfrak{N}$.
\end{theorem}

\begin{proof} 
By Theorem~34 and Remark~32 in \cite{kriz}, our hypothesis on $E[p]$ implies that there is a congruence
\begin{equation}\label{eq:cong-mf}
\source{anticyclotomic-iwasawa-rational-elliptic-eisenstein}
f\equiv G^{}\pmod{p},
\end{equation}
where $G^{}$ is a certain weight two Eisenstein series (denoted $E_{2}^{\phi,\phi^{-1},(N)}$ in \emph{loc.\,cit.}). 
Viewed as a $p$-adic modular form, $G$ defines $\Z_p^{\rm ur}$-valued measures $\mu_{G,\mathfrak{a}}$
on $\Z_p$ by the rule (\ref{eq:def-meas}). With the notations introduced in the proof of Theorem~\ref{thm:BDP}, set
\begin{equation}\label{eq:Eis-measure}
\source{anticyclotomic-iwasawa-rational-elliptic-eisenstein}
\mathscr{L}_{v,\eta}(G,\phi)=
\sum_{[\mathfrak{a}]\in\mathrm{Pic}(\mathcal{O}_c)}
\eta(\mathfrak{a})\mathbf{N}(\mathfrak{a})^{-1}
\int_{\Z_p^\times}\eta_v(\phi\vert[\mathfrak{a}])\;\mathrm{d}\mu_{G^\flat,{\mathfrak{a}}}
\end{equation}
where the cusp form $f$ in $(\ref{eq:def-meas})$ has been replaced by $G$, and let $\mathcal{L}_G^{}\in\Lambda^{\rm ur}$ be the element defined by 
\[
\Lcal_G(\xi):=\mathscr{L}_{v,\eta}(\eta^{-1}\xi).
\]
Then for $\xi$ an arbitrary character of $\Gamma$ crystalline at both $v$ and $\bar{v}$ and corresponding to a Hecke character of $K$ of infinity type $(n,-n)$ for some $n\in\Z_{>0}$ with $n\equiv 0\pmod{p-1}$, the calculation in \cite[Prop.~37]{kriz} (taking $\chi^{-1}=\xi\mathbf{N}_K$, $\psi_1=\phi$, and $\psi_2=\phi^{-1}=\psi\omega^{-1}$ in the notations of \emph{loc.cit.}, so in particular $j=n-1$) shows that
\begin{equation}\label{eq:kriz-prop37}
\source{anticyclotomic-iwasawa-rational-elliptic-eisenstein}
\Lcal_G(\xi)=\frac{\Omega_p^{2n}}{\Omega_\infty^{2n}}\cdot\frac{\Gamma(n+1)\phi^{-1}(-\sqrt{D_K})\xi(\bar{\mathfrak{t}})t}{\mathfrak{g}(\phi)}\cdot\frac{(2\pi i)^{n-1}}{\sqrt{D_K}^{n-1}}\cdot
\Xi_{\xi^{-1}\mathbf{N}_K^{-1}}(\phi,\psi\omega^{-1},N_+,N_-,N_0)\cdot L(\phi_K\xi\mathbf{N}_K,0),
\end{equation}
where $\phi_K$ denotes the base change of $\phi$ to $K$, and
\begin{align*}
\Xi_{\xi^{-1}\mathbf{N}_K^{-1}}(\phi,\psi\omega^{-1},N_+,N_-,N_0)=\prod_{\ell\vert N_+}&(1-\phi^{-1}\xi(\bar{w}))\cdot\prod_{\ell\vert N_-}(1-\phi\xi(\bar{w})\ell^{-1})\\
&\times\prod_{\ell\vert N_0}(1-\phi^{-1}\xi(\bar{w}))(1-\phi\xi(\bar{w})\ell^{-1}).
\end{align*}
Comparing with the interpolation property of $\mathcal{L}_\phi$ in Theorem~\ref{thm:kriz}, and noting that
\[
\mathcal{E}_{\phi,\psi}(\xi^{-1})=\Xi_{\xi^{-1}\mathbf{N}_K^{-1}}(\phi,\psi\omega^{-1},N_+,N_-,N_0)
\]
for all $\xi$ as above, the equality $(\ref{eq:kriz-prop37})$ implies that
\begin{equation}\label{eq:cor-37}
\source{anticyclotomic-iwasawa-rational-elliptic-eisenstein}
\Lcal_G = \mathcal{E}_{\phi,\psi}^\iota\cdot\Lcal_\phi.
\end{equation}

On the other hand, the congruence $(\ref{eq:cong-mf})$ implies the congruences
\[
\binom{\theta}{m}f(x_{\mathfrak{a}})\equiv
\binom{\theta}{m}G(x_{\mathfrak{a}})\pmod{p\Z_p^{\rm ur}}
\]
for all $m\geqslant 0$, which in turn yield the congruence
\begin{equation}\label{eq:cor-meas}
\source{anticyclotomic-iwasawa-rational-elliptic-eisenstein}
\Lcal_E\equiv(\Lcal_G)^2\pmod{p\Lambda^{\rm ur}}.
\end{equation}
The combination of (\ref{eq:cor-37}) and (\ref{eq:cor-meas}) now yields the theorem.
\end{proof}


\begin{theorem}
\source{anticyclotomic-iwasawa-rational-elliptic-eisenstein}
\notready\label{cor:Kriz}
\source{anticyclotomic-iwasawa-rational-elliptic-eisenstein}
Assume that $E[p]^{ss}=\mathbb{F}_p(\phi)\oplus\mathbb{F}_p(\psi)$ as $G_\bQ$-modules, with the characters $\phi, \psi$ labeled so that $p\nmid\mathrm{cond}(\phi)$, and suppose $\phi\neq\mathds{1}$. Then $\mu(\Lcal_E)=0$ and 
\[
\lambda(\Lcal_E)=\lambda(\Lcal_\phi)+\lambda(\Lcal_\psi)+\sum_{w\in S}\bigl\{\lambda(\mathcal{P}_w(\phi))+\lambda(\mathcal{P}_w(\psi))-\lambda(\mathcal{P}_w(E))\bigr\}.
\]
\end{theorem}

\begin{proof}
Since $K$ satisfies (\ref{eq:intro-Heeg}) and (\ref{eq:intro-spl}), the conductors of both $\phi$ and $\psi$ are only divisible by primes split in $K$, and hence the vanishing of $\mu(\mathcal{L}_E)$ follows immediately from the congruence of Theorem~\ref{thm:kriz} and Hida's result \cite{hidamu=0} (note that the factors $\mathcal{P}_w(\phi)$ and $\mathcal{P}_w(\psi)$ also have vanishing $\mu$-invariant, since again the primes $w$ are split in $K$). 

As for the equality between $\lambda$-invariants, note that the involution of $\Lambda$ given by $\gamma\mapsto\gamma^{-1}$ for $\gamma\in\Gamma$
preserves $\lambda$-invariants, and so
\[
\lambda(\mathcal{P}_w(\theta)^2)=\lambda(\mathcal{P}_w(\theta))+\lambda(\mathcal{P}_{\bar{w}}(\theta)),
\]
using that complex conjugation acts as inversion on $\Gamma$. For the term $\mathcal{E}_{\phi,\psi}$ in Theorem~\ref{thm:kriz} we thus have
\[
\lambda((\mathcal{E}_{\phi,\psi}^\iota)^2)=
\lambda((\mathcal{E}_{\phi,\psi})^2)=\sum_{w\vert N_0N_-}\lambda(\mathcal{P}_w(\phi))+\sum_{w\vert N_0N_+}\lambda(\mathcal{P}_w(\psi)),
\]
where $w$ runs over all divisors, not just the ones dividing $\mathfrak{N}$. Using the congruence relations in Theorem~\ref{thm:kriz} (in particular, that $a_\ell\equiv 0\pmod{p}$ for $\ell\vert N_0$) this can be rewritten as
\begin{equation}\label{eq:Euler-comp}
\source{anticyclotomic-iwasawa-rational-elliptic-eisenstein}
\lambda((\mathcal{E}_{\phi,\psi}^\iota)^2)=\sum_{w\in S}\bigl\{\lambda(\mathcal{P}_w(\phi))+\lambda(\mathcal{P}_w(\psi))-\lambda(\mathcal{P}_w(E))\bigr\}.
\end{equation}
On the other hand, since $\psi=\phi^{-1}\omega$, the functional equation for the Katz $p$-adic $L$-function (see e.g. \cite[Thm.~27]{kriz}) yields
\begin{equation}\label{eq:functional}
\source{anticyclotomic-iwasawa-rational-elliptic-eisenstein}
\lambda(\mathcal{L}_\psi)=\lambda(\mathcal{L}_\phi).
\end{equation}
The result now follows from Theorem~\ref{thm:kriz} combined with $(\ref{eq:Euler-comp})$ and $(\ref{eq:functional})$.
\end{proof}

Together with the main result of $\S\ref{sec:algebraic}$, we arrive at the following.

\begin{theorem}
\source{anticyclotomic-iwasawa-rational-elliptic-eisenstein}
\notready\label{mulambda} 
\source{anticyclotomic-iwasawa-rational-elliptic-eisenstein}
Assume that $E[p]^{ss}=\mathbb{F}_p(\phi)\oplus\mathbb{F}_p(\psi)$ with $\phi\vert_{G_p}\neq\mathds{1},\omega$. Then $\mu(\mathcal{L}_E)=\mu(\mathfrak{X}_E)=0$ and 
\[
\lambda(\Lcal_E)=\lambda(\mathfrak{X}_E).
\]
\end{theorem}

\begin{proof}
The vanishing of $\mu(\mathfrak{X}_E)$ (resp. $\mu(\mathcal{L}_E)$) has been shown in Proposition~\ref{propmodp} (resp. Theorem~\ref{cor:Kriz}). On the other hand, 
Iwasawa's main conjecture for $K$ (a theorem of Rubin \cite{rubinmainconj}) yields in particular the equalities
$\lambda(\Lcal_\phi)=\lambda(\mathfrak{X}_\phi)$ and  
$\lambda(\Lcal_\psi)=\lambda(\mathfrak{X}_\psi)$. The combination of Theorem~\ref{MAINalgside} and Theorem~\ref{cor:Kriz} therefore yields the result.
\end{proof}
